\documentclass[12pt]{article}
\usepackage[utf8]{inputenc}
\usepackage[spanish]{babel}
\usepackage{geometry}
\geometry{a4paper, margin=1in}

\title{Perfil Académico y Profesional de la Ingeniería en Mecatrónica}
\author{}
\date{}

\begin{document}

\maketitle

\section{Materias y Áreas de Estudio}

La formación en Ingeniería en Mecatrónica es de carácter multidisciplinario, integrando conocimientos de mecánica, electrónica, informática y control para el diseño de soluciones tecnológicas integrales [1]. Los estudiantes deben poseer facilidad e interés en las \textbf{matemáticas y la física aplicadas} [1].

De acuerdo con el mapa curricular y las áreas principales de formación, las materias se dividen en los siguientes ejes temáticos:

\begin{itemize}
    \item \textbf{Mecánica:} Incluye el diseño de piezas y estructuras, análisis de materiales, sistemas de movimiento, termodinámica y manufactura asistida por computadora (CAD/CAM) [2, 3].
    \item \textbf{Electrónica:} Se estudian circuitos eléctricos, microcontroladores, electrónica digital, sistemas de potencia, instrumentación industrial y sensores [4, 5].
    \item \textbf{Informática y Programación:} Abarca la programación de sistemas, inteligencia artificial, interfaces hombre-máquina (HW-SW), redes industriales, bases de datos y simulación computacional [4, 5].
    \item \textbf{Control y Automatización:} Materias fundamentales como ingeniería de control (I y II), control lógico programable (PLC), robótica (A y B) y diseño de sistemas automatizados [3-5].
\end{itemize}

\section{Campo Laboral}

El egresado en esta carrera tiene la capacidad de identificar, formular y resolver problemas de ingeniería complejos [6]. Podrá desempeñarse en cualquier sector que involucre el diseño, desarrollo e implementación de sistemas automatizados [7].

A continuación, se describen los principales ámbitos de desempeño laboral:

\subsection{Sectores Industriales y Tecnológicos}
\begin{itemize}
    \item \textbf{Industria Tecnológica:} Desarrollo de tecnología para electrónica de consumo, como smartphones, redes y dispositivos para el hogar inteligente (ej. Samsung) [3].
    \item \textbf{Empresas Automotrices:} Diseño de sistemas de control, vehículos eléctricos, autónomos y trenes de aterrizaje (ej. GM, Honda, Cummins) [3, 5].
    \item \textbf{Robótica y Aeroespacial:} Creación de drones, sistemas de navegación, satélites y robots de exploración espacial (ej. NASA/JPL) [3, 5].
    \item \textbf{Automatización Industrial:} Optimización de líneas de producción en industrias metal-mecánica, química, petroquímica y de alimentos [5, 7].
    \item \textbf{Desarrollo de Tecnología Médica:} Diseño de prótesis inteligentes y equipos biomédicos de alta precisión [3, 5].
\end{itemize}

\subsection{Funciones Profesionales}
El ingeniero mecatrónico puede realizar actividades de \textbf{diseño, consultoría, instalación, mantenimiento y dirección de proyectos} [7]. También tiene cabida en el sector académico a través de la \textbf{docencia e investigación} en nuevas tecnologías [7]. Su perfil es altamente demandado por su capacidad para supervisar proyectos de innovación en contextos globales que consideran el impacto económico, ambiental y social [2, 6].

\end{document}